\section[Cubic Subcalibration Criteria]{Cubic Subcalibration Criteria and Radial Models}
\label{sec:cubic-potentials}

The radial table of \cref{sec:cubic-algebraic-theory} gives a degree-six
subcalibration from two invariant inequalities or one sufficient trace
bound. An additional identity yields a degree-seven polynomial alternative.

\subsection{Existence from the radial cubic table}
\label{sec:cubic-unified-sextic}

The common defect formula also applies to the nonradial models of
\cref{sec:nonradial-examples}, using their five-generator table.

\begin{lemma}[The common harmonic potential]
\label{lem:cubic-harmonic-common-potential}
Let \(F\not\equiv0\) be a homogeneous harmonic cubic on \(\R^N\),
\(N\ge2\), satisfying \(\cL(F)=-qF\) for a homogeneous quadratic polynomial
\(q\). Put \(K=\Delta S\) and \(T=\Gamma(S,S)\). For
\(G=F\sqrt{RS}\), there is a homogeneous polynomial
\(\mathcal P_F\) of degree eighteen such that
\begin{equation}\label{eq:cubic-harmonic-common-defect}
 \cL(G)=\frac{F\mathcal P_F}{16\sqrt R\,S^{3/2}}
 \qquad\text{on }\{S>0\}.
\end{equation}
If \(\mathcal P_F\ge0\) there and
\(Z=\{S=0\}\) has locally zero \(\cH^{N-1}\)-measure, then
\(G\) generates a two-sided subcalibration. The phase \(\{F<0\}\)
is locally perimeter minimizing and its boundary current is locally
mass minimizing, with support \(\{F=0\}\).
\end{lemma}

\begin{proof}
Using \(\Gamma(F,S)=2qF\) and homogeneity gives, on \(S>0\),
\begin{align}
 Q_G:=|\nabla G|^2
 &=RS^2+F^2\left[R\left(2q+\frac{T}{4S}\right)+11S\right],
 \label{eq:cubic-harmonic-common-gradient}\\
 \Delta G
 &=G\left[\frac{N+9}{R}+\frac{2q+K/2}{S}
                         -\frac{T}{4S^2}\right].
 \label{eq:cubic-harmonic-common-laplacian}
\end{align}
Substitution into \(\cL(G)=Q_G\Delta G-\Gamma(G,Q_G)/2\)
gives \eqref{eq:cubic-harmonic-common-defect} with a polynomial
\(\mathcal P_F\). Its explicit coefficient formula is
\(\mathcal P_F=C_0+FC_1+F^2C_2\), with
\eqref{eq:cubic-nonharmonic-coefficients} specialized to
\(t=0\) and \(H_2=|\Hess F|^2=K/2\).
\Cref{app:nonharmonic-cubic-certificate} gives the full expansion
and derivation.
Since \(Z\subset\{F=0\}\) by Euler's identity, \(G\) extends
continuously with the signs of \(F\). Its positive multiplier and degree
six give \(\nabla G\ne0\) off \(Z\) by the homogeneous-multiplier
verification following \cref{cor:gradient-subcalibration}.
That corollary and \cref{thm:removable-subcalibration} give comparison;
harmonicity identifies the full support by
\cref{rem:polynomial-finite-perimeter}.
\end{proof}

By \cref{prop:cubic-weight-in-minimal-algebra} and closure,
\(\mathcal P_F\in\cA_R(F)\); its sign is tested on the actual generator
image. No separate condition \eqref{eq:cubic-nonharmonic-PSD} is needed.

For the radial existence criterion, use \(D\) from
\eqref{eq:cubic-unified-D} and put \(K=\Delta S=2\tr(L_x^2)\).
By \cref{lem:cubic-radial-transports}, the criterion depends only
on the radial table and \(K,R,S,D\).

\begin{theorem}[Subcalibrations from the radial cubic table]
\label{thm:cubic-unified-sextic}
Let \(F\not\equiv0\) be a homogeneous harmonic cubic on a Euclidean
space of dimension \(N\ge8\), with \(\Gamma(F,S)=4RF\). Define \(D\)
by \eqref{eq:cubic-unified-D} and put \(K=\Delta S\).
Assume, at every point,
\begin{equation}\label{eq:cubic-unified-trace-bound}
 K\ge16R
\end{equation}
and
\begin{equation}\label{eq:cubic-unified-H}
 \mathscr H:=(N-12)S^2+\frac12R(K-16R)S-12RD\ge0.
\end{equation}
Then the phase \(\{F<0\}\) is locally perimeter minimizing, and its
oriented boundary current is strictly area minimizing with support
\(\{F=0\}\). These conclusions are realized by the degree-six potential
\[
 G=F\sqrt{RS}
   =\frac1{12}|x|\,|x^2|\langle x^2,x\rangle
\]
which is continuous on \(V\), smooth on \(V\setminus Z\), and has
\(\{G<0\}=\{F<0\}\), where \(Z=\{S=0\}\).
Its gradient is nonzero off \(Z\), and its normalized gradient is a
two-sided subcalibration whose divergence has the strict sign of \(F\)
when \(F\ne0\).
\end{theorem}

The polynomial \(\mathscr H\) has a direct boundary interpretation.
Apply \cref{prop:harmonic-power-Jacobi} with \(e=3\),
\(\sigma=1\), \(\gamma=1/2\), \(q=2R\), and
\(T=4RS+12D\). Since \(w=1/(rS)\), this gives
\begin{equation}\label{eq:cubic-H-Jacobi-meaning}
 \left.\mathscr H\right|_{C_{\reg}}
 =RS^2\left(-\frac{J_Cw}{w}\right).
\end{equation}
Since \(G\) has degree six, \cref{thm:ground-state-gap} also gives
\[
 Q_C(u)\geq\left(6-\frac{N-1}{2}\right)^2\int_Cr^{-2}u^2,
 \qquad u\in C_c^\infty(C_{\reg}\setminus\{0\}).
\]
This proves strict stability whenever \(N\ne13\), including when
the link is singular.

\begin{lemma}[Bounds for radial cubics]
\label{lem:cubic-unified-bounds}
Assume the structural hypotheses of \cref{thm:cubic-unified-sextic},
without its dimension and inequality assumptions. Then
\begin{equation}\label{eq:cubic-unified-bounds}
 D\ge0,\qquad S\le\frac{R^2}{3},\qquad
 3RD\le S^2,\qquad 9D\le RS,\qquad F^2\le D.
\end{equation}
\end{lemma}

\begin{proof}
By \cref{lem:cubic-radial-transports}, the last identity in
\eqref{eq:cubic-unified-table} gives \(D\ge0\) on \(S>0\).
At \(S=0\), the nonnegative function \(S\) has zero gradient, and \eqref{eq:cubic-unified-D} gives \(D=0\).
Take a positive maximum of \(S\) on \(R=1\). Homogeneity gives
\(\nabla S=4Sx\) there. Pairing with \(\nabla D\) and using
Euler's identity gives \(24SD=8D\). Thus \(S=1/3\) if \(D>0\).
If \(D=0\), then \(16S^2=\Gamma(S,S)=4S\), so \(S=1/4\).
Homogeneity proves the bound for \(S\).

Cauchy--Schwarz gives \((8RD)^2\le(4RS+12D)(4SD)\).
When \(D>0\), this is
\(D(4R^2-3S)\le RS^2\). Since \(3S\le R^2\), it implies
\(3RD\le S^2\). The assertion is immediate when \(D=0\).
Similarly, Euler's identity and \eqref{eq:cubic-unified-table} give
\[
 36D^2=\langle x,\nabla D\rangle^2\le4RSD,
 \qquad
 4S^2F^2=\Gamma(F,D)^2\le4S^2D.
\]
These prove the final two bounds wherever cancellation is possible.
At \(S=0\), Euler's identity for \(F\) gives \(F=0\), so the
remaining zero cases also follow.
\end{proof}

\begin{proof}[Proof of \cref{thm:cubic-unified-sextic}]
Specializing \cref{lem:cubic-harmonic-common-potential} to
\(q=2R\) and \(T=4RS+12D\) gives, on \(S>0\),
\begin{align}
 Q:=|\nabla G|^2
 &=RS^2+F^2\left(\frac{3DR}{S}+5R^2+11S\right),
 \label{eq:cubic-unified-gradient}\\
 \Delta G
 &=G\left(\frac{N+9}{R}+\frac{K}{2S}
                     +\frac{3R}{S}-\frac{3D}{S^2}\right).
 \label{eq:cubic-unified-laplacian}
\end{align}
In particular, \(Q\ge RS^2>0\). Substituting these formulas into
\(\cL(G)=Q\Delta G-\Gamma(G,Q)/2\) yields
\begin{equation}\label{eq:cubic-unified-defect}
 \cL(G)=\frac{G}{RS^2}
             \left(RS^2\mathscr H+F^2\mathscr J\right),
\end{equation}
where
\begin{align}
 \mathscr J={}&RD\left[\tfrac32R(K-16R)
                  +15(R^2-3S)+3(N-8)S\right]\notag\\
 &+\tfrac52(K-16R)R^3S+45R^4S
             +\tfrac{11}{2}(K-16R)RS^2\notag\\
 &+(5N+33)R^2S^2+(11N+44)S^3.
 \label{eq:cubic-unified-J}
\end{align}
This calculation uses only \cref{lem:cubic-radial-transports} and
homogeneity; no derivative of \(K\) is involved. Every term in \(\mathscr J\)
is nonnegative by \cref{lem:cubic-unified-bounds},
\eqref{eq:cubic-unified-trace-bound}, and \(N\ge8\). The term
\(45R^4S\) is positive. Hence \(F\cL(G)>0\) wherever \(F\ne0\).

The critical set is removable by \cref{lem:cubic-trace-removability}.
Since \(\mathcal P_F=16(RS^2\mathscr H+F^2\mathscr J)\),
\cref{lem:cubic-harmonic-common-potential} gives the minimizing
and support conclusions.
Both open phases are nonempty because \(F\) is odd, and the strict
divergence sign makes both integrals in \eqref{eq:two-phase-flux}
positive. The conical sets \(E=\{F<0\},Z\) and zero-homogeneous field \(X_G\)
therefore satisfy \cref{thm:strict-minimality-flux}, proving strict
area minimality.
\end{proof}

\begin{corollary}[Existence from a trace bound]
\label{cor:cubic-unified-trace}
Under the structural hypotheses of \cref{thm:cubic-unified-sextic},
with \(N\ge8\), the single bound
\begin{equation}\label{eq:cubic-unified-simple-trace}
 \tr(L_x^2)\ge
       \max\left\{8,\frac{40-N}{3}\right\}|x|^2
\end{equation}
implies its two inequality assumptions and all its conclusions.
\end{corollary}

\begin{proof}
The bound implies \(K\ge16R\). By
\cref{lem:cubic-unified-bounds},
\[
 \mathscr H\ge(N-16)S^2+\tfrac12R(K-16R)S.
\]
For \(N\ge16\) this is nonnegative. For \(8\le N<16\), the bound
\(S\le R^2/3\), with the negative coefficient \(N-16\), gives
\[
 \mathscr H\ge RS\left[\frac K2+\frac{N-40}{3}R\right]\ge0.
\]
\end{proof}

The trace test allows nonscalar \(K\), without a closure condition
on \(\Delta D\). The low-dimensional Clifford endpoint instead uses
\eqref{eq:cubic-unified-H} directly.

\subsection{Polynomial alternatives of degree seven}
\label{sec:cubic-algebraic-existence}

To obtain polynomial potentials in place of \(F\sqrt{RS}\), we use
the degree-seven potential \(SF\). Its defect identity holds on the full
radial cubic table. We first impose the additional relation \(D=F^2\);
\cref{prop:exc-SF-certificate} instead assumes a scalar trace coefficient
and allows an independent \(D\).

\begin{lemma}[The gradient-square defect]
\label{lem:cubic-SF-defect}
Let \(F\) be a nonzero homogeneous harmonic cubic satisfying
\(\Gamma(F,S)=4RF\), put \(K=\Delta S\), and define \(D\) by
\eqref{eq:cubic-unified-D}.
For \(G=SF\), one has
\begin{equation}\label{eq:cubic-SF-general-gradient}
 Q:=|\nabla G|^2=S^3+12F^2(RS+D),\qquad
 \Delta G=(K+8R)F,
\end{equation}
and
\begin{equation}\label{eq:cubic-SF-general-defect}
 \begin{aligned}
 \frac{\cL(G)}F={}&(K-16R)S^3-30DS^2\\
   &+F^2\bigl[(12K-72R)D+12RKS-96S^2\bigr].
 \end{aligned}
\end{equation}
The quotient denotes its polynomial extension across \(F=0\).
No scalar-trace condition or identity for \(\Delta D\) is required.
\end{lemma}

\begin{proof}
Apply the product rule and \cref{lem:cubic-radial-transports} to
\(Q=|\nabla(SF)|^2\), then substitute in
\(\cL(SF)=Q\Delta(SF)-\Gamma(SF,Q)/2\).
The individual pairings are recorded in \cref{app:application-details}.
\end{proof}

\begin{corollary}[A polynomial gradient-square certificate]
\label{thm:cubic-algebraic-trace-certificate}
Let \(V\) be a Euclidean commutative metrized algebra of dimension
\(N\ge3\). Suppose that, for every $x\in V$,
\begin{align}
 \tr L_x&=0,\label{eq:cubic-algebraic-trace-zero}\\
 \langle x^2,x^3\rangle
 &=\frac43|x|^2\langle x^2,x\rangle,
 \label{eq:cubic-algebraic-radial-identity}\\
 |x^3|^2&=|x|^2|x^2|^2+
                  \frac13\langle x^2,x\rangle^2,
 \label{eq:cubic-algebraic-norm-identity}\\
 \tr(L_x^2)&\ge8|x|^2.
 \label{eq:cubic-algebraic-trace-bound}
\end{align}
Set
\[
 F(x)=\frac16\langle x^2,x\rangle,\qquad
 R=|x|^2,\qquad S=|\nabla F|^2,\qquad G=SF.
\]
Then $G$ is a polynomial potential of degree seven with
\[
 \{G<0\}=\{F<0\},\qquad
 \{\nabla G=0\}=\{\nabla F=0\}.
\]
Its normalized gradient is a two-sided subcalibration away from
this common critical set, which has locally zero
$\cH^{N-1}$-measure. Hence
$\partial\llbracket\{F<0\}\rrbracket$ is locally mass minimizing.
The divergence has a strict sign at every point with $F\ne0$.
\end{corollary}

\begin{proof}
The multiplication dictionary turns the three identities into
\(\Delta F=0\), \(\Gamma(F,S)=4RF\), and
\(\Gamma(S,S)=4RS+12F^2\), so \(D=F^2\).
Bochner's identity gives \(K=\Delta S=2\tr(L_x^2)\ge16R\),
and \cref{lem:cubic-unified-bounds} gives \(S\le R^2/3\).
Substituting these facts in \cref{lem:cubic-SF-defect} yields the
nonnegative decomposition
\eqref{eq:cubic-variable-trace-positive-decomposition}, which is
strict wherever \(F\ne0\). The same gradient identity identifies
the exceptional set with \(\{\nabla F=0\}\), and
\cref{lem:cubic-trace-removability,cor:gradient-subcalibration}
give the minimizing conclusion. The coefficient decomposition is proved in
\cref{app:application-details}.
\end{proof}

\begin{remark}[Relation to gradient-square closure]
\label{rem:cubic-algebraic-certificate-scope}
The norm identity \eqref{eq:cubic-algebraic-norm-identity} is exactly
\(D=F^2\). In the three-generator closure of
\cref{sec:gradient-square-rank-two},
\cref{prop:cubic-singular-normalization} gives \(K=4kR\), so the
trace bound becomes \(k\ge4\). The quaternionic rank-four trace
model has \(k=5\) and satisfies these hypotheses
(\cref{cor:quaternionic-trace-cubic}). Unlike this certificate,
\cref{thm:cubic-unified-sextic} allows an independent \(D\);
\(SF\) provides a polynomial alternative when the displayed defect
is nonnegative.
\end{remark}

The cubic--algebra dictionary and the multiplier obstruction do not
require harmonicity. For a general quadratic weight in
\(\cL(F)=-qF\), \cref{app:nonharmonic-cubic-certificate} retains
the trace terms in the defect of \(F\sqrt{RS}\).
\Cref{thm:cubic-nonharmonic-certificate} gives a sufficient
polynomial positivity test.

\paragraph{Applications to the classical cubic models.}
The normal forms in \cref{sec:cubic-history} give three-generator
rank-four trace cubics and four-generator complexified rank-three Jordan
norms and \(G_2\) trilinear cubics. The table collects their
\(F\sqrt{RS}\) data and the nonradial model of
\cref{sec:nonradial-examples}.

For the Cartan--M\"unzner cubics \(\Phi\) in dimensions \(14,26\),
set \(F=\Phi/(3\sqrt3)\). Then \(S=R^2/3\), \(D=R^3/27\),
and \(K=4(N+2)R/3\). The common potential is proportional to
\(R^{3/2}F\); the polynomial alternative \(RF\) is given in
\cref{thm:iso-polynomial-table,ex:harmonic-isoparametric-radial}.
\[
\begin{array}{l|r|c}
\text{model}&N&\text{verification data}\\ \hline
\text{Cartan, quaternionic}&14&K/R=64/3\\
\text{Cartan, octonionic}&26&K/R=112/3\\
\text{rank-four quaternionic trace}&27&K/R=20\\
\text{complexified rank-three Jordan norm},\ d=2&18&K/R=16\\
\text{complexified rank-three Jordan norm},\ d=4&30&K/R=24\\
\text{complexified rank-three Jordan norm},\ d=8&54&K/R=40\\
G_2\text{ trilinear}&21&K/R=24\\
\text{nonradial quaternionic model}&26&\text{five-generator closure}
\end{array}
\]
The first seven rows satisfy the radial trace bound
\(K/R\geq\max\{16,2(40-N)/3\}\).
The final row is verified in \cref{thm:nonradial-minimizing}.
In every row the boundary current is strictly area minimizing.

\subsection{Three-generator closure and rank-four trace cubics}
\label{sec:gradient-square-rank-two}

For the closure and realizations in
\cref{prop:cubic-singular-normalization,prop:cubic-gradient-classification},
we determine area minimization, trace models, and the sharp stability
constant. The degree-seven polynomial $G=SF$ complements the common
degree-six potential $F\sqrt{RS}$.

\subsubsection{The tangent obstruction and the existence theorem}

The following closure criterion retains its displayed generality:
for example, the reducible polynomial \(F=xyz\) on \(\R^3\)
satisfies \eqref{eq:cubic-normalized-closure} with \(k=1\), by
direct differentiation. This comparison endpoint is not subject to
the introductory irreducibility convention.

\begin{theorem}[Area minimization from gradient-square closure]
\label{thm:cubic-rank-two-existence}
Let \(F\) be a nonzero homogeneous cubic on \(\R^N\), \(N\geq3\),
and let \(S=|\nabla F|^2\).  Suppose that
\(\R[R,F,S]\) is a Laplacian algebra and that \(F\) has a
nonzero critical point.  Write \(N=6k-3\) as in
\cref{prop:cubic-singular-normalization}.  Then
\begin{equation}\label{eq:cubic-minimality-threshold}
 \partial\llbracket\{F<0\}\rrbracket\text{ is locally mass minimizing}
 \quad\Longleftrightarrow\quad k\geq4.
\end{equation}
Every minimizing member is strictly area minimizing in the sense of
\cref{def:four-variational-notions}, with polynomial potential
\(G=SF\) of degree seven. It is also strictly stable, with the
optimal Hardy constant determined in \cref{thm:sharp}.
\end{theorem}

\begin{proof}
Normalize as in \eqref{eq:cubic-normalized-closure} and suppose
\(k\geq4\). The dictionary in
\cref{prop:cubic-multiplication-dictionary} translates the normalized
gradient identities into the product identities of
\cref{thm:cubic-algebraic-trace-certificate}. Moreover, Bochner's
identity gives
\[
 K=\Delta S=4kR,\qquad
 \tr(L_x^2)=\tfrac12K=2kR\geq8R.
\]
That theorem therefore gives the polynomial subcalibration potential
\(G=SF\), including the required sign and the removability of its
critical set. Its divergence is strictly signed wherever \(F\ne0\).
Since \(G\) is homogeneous and polynomial and both phases are nonempty,
\cref{cor:polynomial-strict-minimality} gives the strict mass gap.

The sharp stability constant is computed in \cref{thm:sharp}.
Harmonicity and \cref{prop:minimizing-irreducibility} also show that
every minimizing member is irreducible.

For the converse, \cref{prop:cubic-gradient-classification} identifies
\(k=1\) with the reducible harmonic cubic \(xyz\), whose phases cannot
minimize perimeter by \cref{prop:minimizing-irreducibility}.
For \(k=2,3\), take a nonzero critical point \(p\).
The Taylor expansion in \cref{prop:cubic-singular-normalization} gives
\[
 r^{-2}F(p+rz)\longrightarrow
 \frac{|p|}{2}\bigl(|u|^2-|v|^2\bigr),
 \qquad z=(u,v,w),\quad u,v\in\R^k.
\]
The transverse quadratic cone is unstable, since its link has
\(|A|^2=2k-2\) and
\[
 \frac{(2k-3)^2}{4}-(2k-2)<0.
\]
By \cref{sec:tangent-instability-transfer}, the Taylor expansion gives
smooth regular graphical convergence. \Cref{lem:tangent-instability-transfer}
therefore makes these cases unstable and excludes area minimization.
\end{proof}

\subsubsection{The rank-four trace realizations}

The trace cubics on $\operatorname{Herm}_4(\mathbb F_d)_0$,
$d=1,2,4$, realize the three exceptional cases in dimensions
$9,15,27$ \cite[Section~6.11.6(iii)]{NadirashviliTkachevVladut14}.
They admit the following uniform calculation.

\begin{corollary}[The rank-four trace cubics]
\label{cor:quaternionic-trace-cubic}
Let \(\mathbb F_d=\R,\mathbb C,\mathbb H\), with real dimension
\(d=1,2,4\), and equip
\[
 V=\operatorname{Herm}_4(\mathbb F_d)_0,\qquad
 \dim_{\R}V=3+6d,
\]
with \(\langle A,B\rangle=\operatorname{Re}\tr(AB)\).
For \(F(A)=\tfrac13\operatorname{Re}\tr(A^3)\),
\(R=\tr(A^2)\), and \(S=|\nabla F|^2\), one has
\[
 \cA_R(F)=\R[R,F,S],\qquad
 \Delta F=0,\quad \Delta S=4(d+1)R,
\]
\[
 \Gamma(F,S)=4RF,\qquad \Gamma(S,S)=4RS+12F^2.
\]
The generators are algebraically independent. The oriented zero cone
is area minimizing if and only if \(d=4\); in that case \(G=SF\)
proves strict area minimality. The common potential \(F\sqrt{RS}\)
also gives a subcalibration in this case.
\end{corollary}

\begin{proof}
Projection onto the trace-free subspace gives
\[
\begin{gathered}
 \nabla F=A^2-\tfrac R4\Id,\qquad
 S=\tr(A^4)-\tfrac{R^2}{4},\\
 (\Hess F)_A(B)=AB+BA-\tfrac12\operatorname{Re}\tr(AB)\Id.
\end{gathered}
\]
Write \(p_j=\sum_i\lambda_i^j\) after diagonalizing \(A\), where
\(p_1=0\). If \(P\) is the orthogonal projection onto the
sum-zero diagonal vectors, the diagonal Hessian is
\(2P\operatorname{diag}(\lambda_i)P\), with trace zero and
squared Hilbert--Schmidt norm \(2R\). Each off-diagonal entry
contributes \(d\) eigenvalues \(\lambda_i+\lambda_j\). Thus
\[
 \Delta F=0,\qquad
 |\Hess F|^2=2R+d\sum_{i<j}(\lambda_i+\lambda_j)^2
             =2(d+1)R,
\]
and \(\Delta S=4(d+1)R\). Also,
\(\nabla S=4A^3-RA-p_3\Id\). Newton's identities give
\[
 p_5=\tfrac56Rp_3,\qquad
 p_6=\tfrac34Rp_4-\tfrac18R^3+\tfrac13p_3^2.
\]
Substitution proves the two gradient identities. Together with
Euler's identities, \cref{prop:MR-generator-test} gives Laplacian
closure. Since \(S=\Gamma(F,F)\) belongs to \(\cA_R(F)\),
this is the smallest algebra containing \(R,F\).
At distinct eigenvalues, the Vandermonde determinant shows that
\(p_2,p_3,p_4\) have independent differentials on \(p_1=0\);
hence \(R,F,S\) are algebraically independent.

The gradient formula and the transverse model in
\cref{prop:cubic-singular-normalization} identify the singular set as
\[
 \Sing\{F=0\}
 =\{A\in V:A^2=\tfrac14\tr(A^2)\Id\}.
\]
Its nonzero matrices have eigenvalues \((\lambda,\lambda,-\lambda,-\lambda)\),
with \(\lambda>0\). In particular,
\(\operatorname{diag}(1,1,-1,-1)\) is a nonzero critical point.
Since \(\Delta S=4kR\) with \(k=d+1\),
\cref{thm:cubic-rank-two-existence} proves all the minimizing and
nonminimizing assertions. For the common potential, put \(D=F^2\).
The displayed identities imply
\[
 \Gamma(F,D)=2SF,\qquad\Gamma(S,D)=8RD,\qquad
 \Gamma(D,D)=4SD.
\]
For \(d=4\), \(N=27\) and \(\tr L_x^2=10R\), so
\cref{cor:cubic-unified-trace} applies.
\end{proof}

\Needspace{12\baselineskip}
\subsubsection{Exact link geometry and the sharp Hardy constant}
On the sphere write \(f=F|_{\Sph^{N-1}}\) and
\(s=S|_{\Sph^{N-1}}\).
Let \(M=Z^*\cap \Sph^{N-1}\). Since \(Z^*\) is conical and radial
vectors are transverse to the sphere, \(M\) is a compact smooth
submanifold of dimension \(4k-4\). The zero link
\(\Sigma=\{F=0\}\cap \Sph^{N-1}\) has regular part of dimension \(6k-5\).
Its transverse dimension along \(M\) is thus \(2k-1\).
If \(\rho\) denotes spherical distance to \(M\), then, near \(M\),
\begin{equation}\label{eq:tube}
 s\asymp\rho^2,\qquad
 \cH^{N-2}\bigl(\Sigma_{\reg}\cap\{\rho<\eps\}\bigr)
 \leq C\eps^{2k-1}.
\end{equation}
The proof of these estimates is given in \cref{sec:cubic-local-geometry}.

\paragraph{Two geometric identities.}
With \(\nu=\nabla F/\sqrt S\), the curvature identity proved in
\cref{prop:harmonic-power-Jacobi} is
\(|A_C|^2=(\Delta S)/(2S)-|\nabla S|^2/(2S^2)\).
Substituting \(\Delta S=4kR\) and \(|\nabla S|^2=4RS\)
on the regular cone gives
\begin{equation}\label{eq:Aexact}
 |A_C|^2=\frac{2(k-1)R}{S},\qquad
 |A_\Sigma|^2=\frac{2(k-1)}s.
\end{equation}

The spherical cross term \(\Gamma_{\Sph}(f,s)=4f(1-3s)\) vanishes on
the zero link. Therefore
\begin{equation}\label{eq:linksgrad}
 |\nabla_\Sigma s|^2=4s-16s^2.
\end{equation}
In particular \(0<s\leq1/4\) on \(\Sigma_{\mathrm{reg}}\), improving
the global bound \(s\leq1/3\) in
\eqref{eq:cubic-gradient-only-S-bound}.

To compute the link Laplacian, differentiate the closure table:
\begin{align*}
 \Gamma(F,\Gamma(F,S))&=4RS+24F^2,\\
 \Hess F(\nabla F,\nabla S)&=\tfrac12|\nabla S|^2=2RS+6F^2.
\end{align*}
Their difference is \(\Hess S(\nabla F,\nabla F)\); hence
\(\Hess S(\nu,\nu)=2R\) on the cone. Since \(S\) has degree four,
its spherical normal Hessian at \(r=1\) equals \(2-4s\).
Subtract this from \(\Delta_{\Sph}s=4k-4(N+2)s\). The link is minimal,
so there is no mean-curvature term, and
\begin{equation}\label{eq:linksLap}
 \Delta_\Sigma s=4k-2-4(N+1)s
 =4k-2-(24k-8)s.
\end{equation}

\paragraph{A positive eigenfunction and its domain.}
Let \(J_\Sigma=\Delta_\Sigma+|A_\Sigma|^2\).
For any real \(p\), equations \eqref{eq:Aexact}--\eqref{eq:linksLap} give
\begin{equation}\label{eq:powerJ}
 J_\Sigma(s^{-p})
 =\left\{\frac{4p^2+(6-4k)p+2k-2}{s}
           +4p(6k-6-4p)\right\}s^{-p}.
\end{equation}
For \(k\geq4\), define
\begin{equation}\label{eq:p}
 \delta_k=\sqrt{4k^2-20k+17}>0,\qquad
 p_k=\frac{2k-3-\delta_k}{4}>0,\qquad
 B_k=4p_k(6k-6-4p_k).
\end{equation}
The coefficient of \(s^{-1}\) in \eqref{eq:powerJ} vanishes, so
\begin{equation}\label{eq:eig}
 \phi=s^{-p_k}>0,\qquad J_\Sigma\phi=B_k\phi.
\end{equation}

\begin{theorem}[Optimal Hardy stability constant]\label{thm:sharp}
Under \eqref{eq:cubic-normalized-closure}, with a nonvertex critical point and \(k\geq4\),
the Friedrichs bottom of \(-J_\Sigma\), initially defined on
\(C_c^\infty(\Sigma_{\mathrm{reg}})\), is \(-B_k\).
The optimal constant in
\[
 Q_C(u)\geq H_k\int_Cr^{-2}u^2,
 \qquad u\in C_c^\infty(C_{\mathrm{reg}}\setminus\{0\}),
\]
is
\begin{equation}\label{eq:sharpH}
 \boxed{H_k=\bigl(k+\sqrt{4k^2-20k+17}\bigr)^2.}
\end{equation}
In particular, for \(k=5\), \(N=27\), \(H_5=42+10\sqrt{17}\).
\end{theorem}

\begin{proof}
For every compactly supported smooth \(v\) on the regular link,
the positive-solution identity applied to \eqref{eq:eig} gives
\begin{equation}\label{eq:linkground}
 \int_\Sigma\bigl(|\nabla v|^2-|A_\Sigma|^2v^2\bigr)
 =\int_\Sigma\phi^2\left|\nabla(v/\phi)\right|^2-B_k\int_\Sigma v^2.
\end{equation}
The bottom is therefore at least \(-B_k\).

By \eqref{eq:tube}, \(\phi\asymp\rho^{-2p_k}\), and \(2k-3-4p_k=\delta_k>0\).
Thus \(\phi\in L^2(\Sigma)\), with finite gradient and potential
energy near \(M\). Choose smooth cutoffs \(\eta_\epsilon\), zero
when \(\rho<\epsilon\), one when \(\rho>2\epsilon\), with
\(|\nabla\eta_\epsilon|\leq C/\epsilon\). The local volume bound gives
\[
 \int_\Sigma\phi^2|\nabla\eta_\epsilon|^2
 =O(\epsilon^{2k-3-4p_k})=O(\epsilon^{\delta_k})\longrightarrow0.
\]
Use \(v=\eta_\epsilon\phi\) in \eqref{eq:linkground}. Its Rayleigh
quotient tends to \(-B_k\), proving the claimed bottom.
The same estimates give convergence to \(\phi\) in the Friedrichs
form norm. Testing \eqref{eq:eig} first against compactly supported
smooth functions and then using form density shows that \(\phi\)
belongs to the Friedrichs operator domain, with eigenvalue \(-B_k\)
for \(-J_\Sigma\).

The cone dimension is \(n=N-1=6k-4\).  By
\eqref{eq:exact-Hardy-gap}, the optimal cone constant is
\[
 \frac{(n-2)^2}{4}+\inf\operatorname{spec}(-J_\Sigma)
 =9(k-1)^2-B_k.
\]
Finally,
\[
 9(k-1)^2-B_k
 =\bigl(3(k-1)-4p_k\bigr)^2=(k+\delta_k)^2,
\]
which proves \eqref{eq:sharpH}.
\end{proof}

\begin{remark}
For the larger root in \eqref{eq:p},
\(2k-3-4p=-\delta_k\), so the cutoff energy does not tend to zero.
The radicand \(\delta_k^2\) equals four times the Hardy stability
constant of the transverse balanced quadratic cone:
\[
 \delta_k^2=4\left[\frac{(2k-3)^2}{4}-(2k-2)\right].
\]
\end{remark}

\subsection{Four-generator radial models}
\label{sec:exceptional-four-generator}

\subsubsection{Complexified Jordan norms}
\label{exc-jordan-models}

The sextic \(D\) is now independent: \(D=(\Gamma(S,S)-4RS)/12\).
In the preceding trace family, \(D=F^2\).


Let \(\mathbb K_d=\R,\mathbb C,\mathbb H,\mathbb O\) for
\(d=1,2,4,8\), respectively, and set
\(J_d=\operatorname{Herm}_3(\mathbb K_d)\), writing \(J=J_d\) locally.
Its Jordan product is \(x\bullet y=(xy+yx)/2\), its unit is \(e\),
and its Euclidean inner product is \(\tr(x\bullet y)\).
Thus \((J_d)_0=e^\perp=\{x\in J_d:\tr x=0\}\) denotes its trace-free part.
On the real vector space
underlying \(V=J\otimes_{\R}\mathbb C\), use the real part of
the Hermitian extension of this inner product. Thus
\(|x+iy|^2=|x|^2+|y|^2\) and \(\dim_{\R}V=6(d+1)\).
Let \(N\) be the complexification
of the Jordan determinant. A Jordan frame consists of three
idempotents \(c_i\) with \(c_i\bullet c_j=\delta_{ij}c_i\),
\(\tr c_i=1\), and \(\sum_i c_i=e\). The normalization is
\[
 N(z_1c_1+z_2c_2+z_3c_3)=z_1z_2z_3,
 \qquad |z_1c_1+z_2c_2+z_3c_3|^2=\sum_i|z_i|^2.
\]
These are the complexified rank-three models in
\cite[Section~6.11.6(iv)]{NadirashviliTkachevVladut14}; the polynomial
there is \(2\operatorname{Re}N\), with the same zero set.
Their REC characteristic dimensions are \((n_1,n_2)=(1,3d+2)\).

The determinant has the coordinate formula
\begin{equation}\label{exc-norm-coordinates}
 z=\begin{pmatrix}
 a&\bar u&\bar v\\ u&b&w\\ v&\bar w&h
 \end{pmatrix},\qquad
 N(z)=abh-a\,n(w)-b\,n(v)-h\,n(u)
       +2\operatorname{Sc}((\bar u w)v).
\end{equation}
Here \(a,b,h\in\mathbb C\) and
\(u,v,w\in\mathbb K_d\otimes_{\R}\mathbb C\);
the bar is algebra conjugation fixing complex scalars,
\(n(u)=u\bar u\), and \(\operatorname{Sc}\) is the complex-linear
extension of the scalar part. Thus \(n\) is a complex quadratic
form, distinct from the Hermitian norm \(|u|^2\).

The quadratic adjoint \(z^\#\) satisfies
\((dN)_z(v)=\tr(z^\#\bullet v)\), with the complex bilinear
trace form. The rank-three identities are
\begin{equation}\label{exc-adjoint}
 z\bullet z^\#=N(z)e,\qquad
 (z^\#)^\#=N(z)z,\qquad
 \tr(z\bullet z^\#)=3N(z).
\end{equation}
They extend from \(J\) to \(V\) as polynomial identities; see
\cite[Section~3.2.7]{NadirashviliTkachevVladut14} for the Jordan algebra
background. Put
\begin{equation}\label{exc-invariants}
 F=\operatorname{Re}N,\qquad R=|z|^2,\qquad
 S=|z^\#|^2,\qquad D=|N(z)|^2.
\end{equation}
Holomorphicity gives \(\Delta F=0\) and \(|\nabla F|^2=S\).
The Jordan product on \(J\) is distinct from the metrized product
obtained by polarizing \(F\) on \(V\).

\begin{theorem}[The smallest closure of the complexified models]
\label{thm:exc-jordan-closure}
For each \(d\in\{1,2,4,8\}\), the four polynomials in
\eqref{exc-invariants} are algebraically independent and
\[
 \cA_R(F)=\R[R,F,S,D].
\]
Their degrees are \(2,3,4,6\); in particular,
\((\cA_R(F))_2=\R R\) and \((\cA_R(F))_3=\R F\).
The Laplacian identities are
\begin{equation}\label{exc-jordan-laplacians}
 \Delta R=12(d+1),\qquad \Delta F=0,\qquad
 \Delta S=4(d+2)R,\qquad \Delta D=4S.
\end{equation}
The gradient pairings, apart from those determined by
\(\Gamma(R,P)=2\deg(P)P\), are
\begin{equation}\label{exc-gradient-table}
\begin{aligned}
 \Gamma(F,F)&=S,& \Gamma(F,S)&=4RF,& \Gamma(F,D)&=2SF,\\
 \Gamma(S,S)&=4RS+12D,& \Gamma(S,D)&=8RD,&
 \Gamma(D,D)&=4SD.
\end{aligned}
\end{equation}
In particular, \(\cL(F)=-2RF\), so the regular zero set is minimal.
Throughout \(V\),
\begin{equation}\label{exc-global-inequalities}
 R^2\geq3S,\qquad RS\geq9D,\qquad S^2\geq3RD,
 \qquad 0\leq F^2\leq D.
\end{equation}
\end{theorem}

\begin{proof}
\emph{Differential identities.}
For real \(a\in J\), \(T_a:z\mapsto a\bullet z\)
is self-adjoint for the Hermitian trace metric, and
\eqref{exc-adjoint} gives \((dN)_z(T_az)=(\tr a)N(z)\).
Consequently \(\exp(itT_a)\) is unitary and multiplies \(N\)
by \(\exp(it\tr a)\). Real Jordan automorphisms are also
unitary and preserve \(N\). These transformations preserve
\(R,S,D\) and transport identities involving \(N\) with its
phase factor.

The linearized Jordan identity implies that the commutators
\([T_a,T_b]\) are derivations, whose exponentials are automorphisms.
Take the standard diagonal frame \((c_i)\). At
\(x=\sum_i r_ic_i\), with distinct positive \(r_i\), the
Peirce decomposition \cite[Section~3.2.6]{NadirashviliTkachevVladut14} is
\[
 J=\bigoplus_i\R c_i\ \oplus\!\bigoplus_{i<j}J_{ij},
 \qquad \dim_{\R}J_{ij}=d.
\]
Here \(J_{ij}\) consists of matrices supported in the \((i,j)\)
and \((j,i)\) entries, and
\(c_i\bullet a=c_j\bullet a=a/2\) for \(a\in J_{ij}\).
Varying the \(r_i\) supplies the real diagonal directions, while
the derivations and unitary infinitesimal transformations give
\[
 [T_{c_i},T_a]x=\tfrac14(r_j-r_i)a,\qquad
 iT_ax=\tfrac12(r_i+r_j)ia,\qquad iT_{c_i}x=ir_ic_i.
\]
These directions span the underlying real space of \(V\).
The map obtained by varying the diagonal entries and the transformation
parameters is therefore a submersion at \(x\), so its image contains
a real open subset of \(V\).

On the diagonal, with \(t_i=r_i^2\),
\(R=e_1(t), S=e_2(t), D=e_3(t)\), where
\(e_1(t)=\sum_i t_i\), \(e_2(t)=\sum_{i<j}t_it_j\), and
\(e_3(t)=t_1t_2t_3\). The first derivatives of
these three invariants in the off-diagonal directions vanish.
Differentiation gives the last row of \eqref{exc-gradient-table}.
Extending \(\Gamma\) complex bilinearly, diagonal differentiation
gives \(\Gamma(N,S)=4RN\), whose real part is the desired identity
for \(F\). Holomorphicity gives the remaining pairings with
\(F\). The complex Hessian of \(N\) has diagonal-block
squared norm \(2R\); on \(J_{ij}\) its block is
\(-r_k\Id\), where \(\{i,j,k\}=\{1,2,3\}\), contributing
\(dR\) in total. Therefore
\[
\begin{gathered}
 |\Hess_{\mathbb C}N|^2=(d+2)R,\qquad
 |\Hess_{\R}F|^2=2(d+2)R,\\
 \Delta S=2|\Hess_{\R}F|^2=4(d+2)R.
\end{gathered}
\]
Here \(|\Hess_{\R}\operatorname{Re}(e^{i\theta}N)|^2
=2|\Hess_{\mathbb C}N|^2\), so the real Hessian norm is
unchanged by the determinant phase.
Also \(\Delta|N|^2=4|z^\#|^2\). Phase covariance transports
these identities to the open subset just obtained, and polynomiality
then proves them on all of \(V\).

\emph{The smallest closure.}
The differential identities satisfy
\cref{prop:cubic-scalar-killing-closure}, which gives the claimed
smallest algebra. At distinct positive \(t_i\), the Jacobian
of \((e_1,e_2,e_3)\) is nonzero. Changing the overall complex
phase preserves \(R,S,D\) and changes \(F\) whenever \(N\)
is nonreal. Hence the four differentials are independent at such
a point.

\emph{Global inequalities.}
The gradient table and \cref{lem:cubic-unified-bounds} give
\eqref{exc-global-inequalities}.
\end{proof}

\subsubsection{Polynomial potentials and the Jordan minimizing range}
\label{exc-multipliers}

The following polynomial alternative also applies to the model in
\cref{exc-g2-model}.

\begin{proposition}[An algebraic sign criterion for the multiplier \(S\)]
\label{prop:exc-SF-certificate}
Let \(F\) be a harmonic cubic on a Euclidean space \(V\), set
\(R=|x|^2\) and \(S=|\nabla F|^2\), and let \(D\) be a
sextic satisfying \eqref{exc-gradient-table} and
\eqref{exc-global-inequalities}. Suppose \(\Delta S=cR\).
For \(G=SF\), one has \(\cL(G)=F\mathcal P_c\), where
\begin{equation}\label{exc-SF-certificate}
\begin{aligned}
 \mathcal P_c={}&(c-\tfrac{58}{3})RS^3
       +\tfrac{10}{3}S^2(RS-9D)\\
 &+F^2\bigl[(12c-72)RD+12cS(R^2-3S)+(36c-96)S^2\bigr].
\end{aligned}
\end{equation}
If \(c\geq58/3\), then \(G\cL(G)\geq0\), and
\(\nabla G\neq0\) wherever \(S>0\).
\end{proposition}

\begin{proof}
Set \(K=cR\) in \eqref{eq:cubic-SF-general-defect} and regroup
to obtain \eqref{exc-SF-certificate}. Every summand is nonnegative
by \eqref{exc-global-inequalities} when \(c\geq58/3\).
The gradient formula \eqref{eq:cubic-SF-general-gradient} gives
\(|\nabla G|^2\geq S^3\).
\end{proof}

\begin{theorem}[Area minimization for the complexified Jordan models]
\label{thm:exc-jordan-minimizing}
For the four models in \cref{exc-jordan-models},
\(\{F<0\}\) is locally perimeter minimizing if and only if
\(d\in\{2,4,8\}\). In all three minimizing dimensions,
\(G=F\sqrt{RS}\) generates a subcalibration. For \(d=4,8\),
the polynomial potential \(SF\) does so as well.
\end{theorem}

\begin{proof}
For \(d=2,4,8\), the dimension and trace coefficient are
\[
 (N,K/R)=(18,16),\ (30,24),\ (54,40),\qquad K=\Delta S.
\]
Since \(\tr L_x^2=K/2\), each pair satisfies
\cref{cor:cubic-unified-trace}; this proves the minimizing assertions
for the common potential. For \(d=4,8\),
\cref{prop:exc-SF-certificate} additionally gives the polynomial
potential \(SF\).

The singular set is the Jordan rank-at-most-one cone
\[
 \Sing\{F=0\}=\Sigma=\{z\in V:z^\#=0\}.
\]
To describe its nonzero stratum, take
\(a\neq0\) in \eqref{exc-norm-coordinates} and set
\[
 B=b-\frac{n(u)}a,\quad C=h-\frac{n(v)}a,\quad
 W=w-\frac{u\bar v}a.
\]
The determinant formula gives
\begin{equation}\label{exc-schur}
 N(z)=a\bigl(BC-n(W)\bigr).
\end{equation}
This identity also holds for \(d=8\). Indeed, multiplicativity
of \(n\) and associativity of the scalar part give
\[
 n\!\left(w-\frac{u\bar v}{a}\right)
 =n(w)-\frac{2}{a}\operatorname{Sc}((\bar u w)v)
       +\frac{n(u)n(v)}{a^2},
\]
which proves \eqref{exc-schur} after expanding \(aBC\).
The coordinate change is holomorphically invertible, so
\(dN=0\) is exactly \(B=C=W=0\) in this chart. Hence
these points are geometrically singular: the transverse real
quadratic form \(\operatorname{Re}(a(BC-n(W)))\) is nondegenerate
and indefinite. Moreover,
\(\Sigma\setminus\{0\}\) is smooth of complex dimension
\(1+2d\), and real codimension \(2(d+2)\). To see that these
charts cover the nonzero set, recall that the spectral theorem makes
primitive idempotents span \(J\), and that real Jordan automorphisms
act transitively on them
\cite[Theorem~III.1.2 and Corollary~IV.2.7]{FarautKoranyi94}.
Some primitive idempotent therefore has nonzero trace pairing with
any nonzero \(z\), and an automorphism takes it to \(c_1\).

For \(d=4,8\), the codimension bound and
\eqref{eq:cubic-SF-general-gradient} allow
\cref{cor:gradient-subcalibration} to be applied to \(SF\).
Harmonicity identifies the boundary-current support with \(\{F=0\}\)
by \cref{rem:polynomial-finite-perimeter}.

For \(d=1\), take \(p\in\Sigma\setminus\{0\}\) and put
\(H=\Hess F(p)\). The chart \eqref{exc-schur} shows that
\(H\) has real rank \(2(d+2)\).
\Cref{prop:cubic-general-singular-spectrum}, with \(q=2R\), gives
\(H^3=R(p)H\). Since \(\tr H=0\), its nonzero eigenvalues
are \(\pm\sqrt{R(p)}\), each of multiplicity \(d+2\).
Thus the tangent cone, with the induced Euclidean metric, is
\begin{equation}\label{exc-jordan-tangent}
 C_{d+1,d+1}\times\R^{4d+2},\qquad
 C_{q,q}=\{(x,y)\in\R^{q+1}\oplus\R^{q+1}:|x|=|y|\}.
\end{equation}
For \(d=1\), the normal factor \(C_{2,2}\) is unstable by the
Hardy criterion in \cref{sec:tangent-instability-transfer}.
\Cref{lem:tangent-instability-transfer} therefore rules out local
perimeter minimization.
\end{proof}

For \(d=2\), complexification identifies \(V\) isometrically
with \(M_3(\mathbb C)\) with its Frobenius metric, and \(N\)
with the determinant. In the notation of \cref{sec:complex-determinants},
\(R=E_1\), \(S=E_2\), and \(D=E_3\). Thus both the closure
table and the potential above are the three-by-three cases of the
construction there.

\subsubsection{The \texorpdfstring{$G_2$}{G2} trilinear model}
\label{exc-g2-model}

Let \(\times\) be the standard cross product on
\(\operatorname{Im}\mathbb O=\R^7\), and set
\[
 F(a,b,c)=\langle a\times b,c\rangle,
 \qquad (a,b,c)\in(\R^7)^3.
\]
Equivalently, \(F(a,b,c)=\varphi(a,b,c)\), where \(\varphi\)
is the associative calibration on \(\R^7\) of Harvey and Lawson
\cite[Section~IV.1.A]{HarveyLawson82a}.
This is the classical model in
\cite[Section~6.4, equation~(6.26), Proposition~6.3]{TkachevClassification10},
whose defining polynomial \(\operatorname{Re}(a(bc))\) is
\(-F\). Up to similarity, its subtracted Jordan model is the cubic
norm restricted to the trace-orthogonal complement of
\(\operatorname{Herm}_3(\R)\) in \(\operatorname{Herm}_3(\mathbb O)\);
see \cite[Proposition~6.4.4 and Table~6.4.1]{NadirashviliTkachevVladut14}.
Let \(A\) be the real \(7\times3\) matrix with
columns \(a,b,c\), and put
\[
 B=A^tA,\qquad R=\tr B,\qquad S=e_2(B),\qquad D=\det B,
\]
where \(e_2(B)\) is the second elementary symmetric polynomial
of the eigenvalues of \(B\).

The compact real group \(G_2=\operatorname{Aut}(\mathbb O)\)
acts diagonally on \(V=(\R^7)^3\). Schwarz's invariant-generation
theorem \cite[Theorem~4 and equations~(5.1)--(5.3)]{Schwarz83} gives
\[
 \R[V]^{G_2}=\R[B_{ij},F].
\]
The commuting \(SO(3)\) action on the columns fixes \(F\)
and conjugates \(B\). Averaging polynomial expressions over
\(SO(3)\), whose conjugation invariants on symmetric matrices
are generated by the elementary symmetric polynomials, therefore gives
\[
 \R[V]^{G_2\times SO(3)}=\R[R,F,S,D].
\]
This invariant ring equals \(\cA_R(F)\):

\begin{theorem}[The trilinear model]
\label{thm:exc-g2}
The four polynomials \(R,F,S,D\) are algebraically independent,
and \(\cA_R(F)=\R[R,F,S,D]\). In particular,
\((\cA_R(F))_2=\R R\) and \((\cA_R(F))_3=\R F\).
They satisfy
\eqref{exc-gradient-table} and \eqref{exc-global-inequalities},
with
\begin{equation}\label{exc-g2-laplacians}
 \Delta R=42,\qquad \Delta F=0,\qquad
 \Delta S=24R,\qquad \Delta D=10S.
\end{equation}
Both \(G=F\sqrt{RS}\) and the polynomial potential \(SF\)
give subcalibrations of \(\{F<0\}\subset\R^{21}\), which is
locally perimeter minimizing.
\end{theorem}

\begin{proof}
The cross-product identity
\(|a\times b|^2=|a|^2|b|^2-\langle a,b\rangle^2\)
gives \(|\nabla F|^2=S\), while trilinearity gives
\(\Delta F=0\). Alternation gives
\(\langle\nabla_{a_i}F,a_j\rangle=\delta_{ij}F\), with
\(a_i\) denoting the columns. In matrix notation,
\[
 \nabla_A S=2A(RI-B),\qquad
 \nabla_A D=2A\operatorname{adj}B.
\]
These formulas and the three-by-three Newton identities give
\eqref{exc-gradient-table}. Direct differentiation of the Gram
invariants gives \(\Delta S=24R\) and \(\Delta D=10S\).
The global inequalities \eqref{exc-global-inequalities} follow
from \cref{lem:cubic-unified-bounds}.

The smallest-algebra assertion now follows from
\cref{prop:cubic-scalar-killing-closure}.
For independence,
choose orthonormal \(e_1,e_2\), put \(e_3=e_1\times e_2\),
and choose a unit \(e_4\) orthogonal to their span. At
\[
 a=r_1e_1,\qquad b=r_2e_2,\qquad
 c=r_3(\cos\theta\,e_3+\sin\theta\,e_4),
\]
the three invariants \(R,S,D\) are the elementary symmetric
functions of \(r_1^2,r_2^2,r_3^2\), while
\(F=r_1r_2r_3\cos\theta\). Distinct positive \(r_i^2\)
and \(0<\theta<\pi\) give four independent differentials.

Here \(N=21\) and \(K=24R\), so
\cref{cor:cubic-unified-trace} gives the common potential.
For the polynomial potential, \cref{prop:exc-SF-certificate} applies
with \(c=24\). The singular set of the zero cone is
\[
 \Sing\{F=0\}=\Sigma=\{S=0\}
 =\{A\in M_{7,3}(\R):\operatorname{rank}A\leq1\}.
\]
At a nonzero rank-one matrix, a column rotation reduces to
\((a,0,0)\), whose transverse leading form
\(\langle a\times\delta b,\delta c\rangle\) is indefinite of rank twelve.
The nonzero stratum has dimension
\(7+3-1=9\), hence codimension \(12\). The gradient bound
\eqref{eq:cubic-SF-general-gradient} and \cref{cor:gradient-subcalibration}
therefore give the same minimizing phase.
\end{proof}

\begin{remark}
For both families, Bochner's identity gives
\(2|L_x|_{\mathrm{HS}}^2=c|x|^2\), where \(c>0\). Hence
\(L_v=0\) implies \(v=0\), and their smallest closures have the intrinsic
degree pattern \((3)\) of \cref{thm:cubic-one-dimensional-degree-three}.
For the \(G_2\) model this differs from its trilinear coordinate
presentation: the larger invariant algebra
\(\R[V]^{G_2}\) contains their separate radii, whereas the
smallest algebra is also invariant under column rotations and
does not contain \(|a|^2,|b|^2,|c|^2\).
\end{remark}



\subsection{Area-minimizing Clifford cubic cones}
\label{sec:Clifford-classification}
\begin{theorem}[Area-minimizing classification]
\label{thm:clifford-classification}
For every admissible symmetric Clifford system with \(q\geq1\), the
oriented multiplicity-one cone \(\mathcal C_{\mathcal A}\) is area
minimizing if and only if
\begin{equation}\label{eq:clifford-classification-list}
 \boxed{m\geq4,\qquad q\geq3,\qquad(m,q)\neq(4,3).}
\end{equation}
Every area-minimizing member in this family is strictly area minimizing.
\end{theorem}

The positive range uses \(G=F\sqrt{RS}\) and
\cref{thm:cubic-unified-sextic}; singular-stratum obstructions exclude
the low-dimensional cases. The remaining pair \((4,3)\) requires a
link variation:

\begin{theorem}[Instability of the pair \((4,3)\)]
\label{thm:Clifford-43-instability}
For every Clifford system with \((m,q)=(4,3)\), the cone
\(\cC_\cA\subset\R^{12}\) is unstable.  More precisely, if
\(\Lambda=\cC_\cA\cap\Sph^{11}\), then
\(\phi=S^{3/4}(1-S/2)\) belongs to the energy space
\(\mathcal D(\Lambda)\) defined in
\cref{sec:instability} and satisfies
\[
 \int_{\Lambda_{\reg}}
 \left(|\nabla_\Lambda\phi|^2-|A_\Lambda|^2\phi^2
       +\frac{81}{4}\phi^2\right)\dd\cH^{10}<0,
\]
where \(A_\Lambda\) is the second fundamental form of
\(\Lambda\subset\Sph^{11}\).
Every tangent cone of \(\cC_\cA\) at a nonzero point is area minimizing.
At such a singular point it is the product of the seven-dimensional Simons
cone and \(\R^4\).
\end{theorem}

\subsection{The Clifford positivity range}
\label{sec:Clifford-common-potential}

\begin{proposition}[The common potential for Clifford cubics]
\label{prop:Clifford-common-potential}
Suppose either \(m\geq8\), \(q\geq3\), or \(m\geq4\), \(q\geq4\).
Then
\[
 G=F\sqrt{RS}
\]
generates a two-sided subcalibration for \(\{F<0\}\), and
\(\cC_{\cA}\) is strictly area minimizing with the prescribed orientation.
The gradient of \(G\) is nonzero on \(\{S>0\}\), and
\(F\cL(G)>0\) wherever \(F\ne0\).
\end{proposition}

\begin{proof}
For \cref{thm:cubic-unified-sextic}, use the normalization
\begin{equation}\label{eq:Clifford-normalized-data}
 \begin{aligned}
 f&=\frac F2,\qquad s=|\nabla f|^2=\frac S4=uv+\frac T4,
 \qquad D=\frac{vT}{4},\\
 \kappa&=\Delta s=\frac K4=4mv+4(q+1)u.
 \end{aligned}
\end{equation}
The closure table \eqref{eq:clifford-closure-table} gives
\begin{equation}\label{eq:Clifford-unified-table}
 \Gamma(f,s)=4Rf,\qquad \Gamma(s,s)=4Rs+12D.
\end{equation}
Thus \(\cL(f)=-2Rf\), and the remaining transport identities are
\eqref{eq:cubic-unified-table}, with \((F,S)\) replaced by \((f,s)\).
If \(m\geq8\) and \(q\geq3\), then
\(N=2m+q+1\geq20\) and \(\kappa\geq16R\).
The trace criterion \cref{cor:cubic-unified-trace} applies directly.

For the second range, write the remaining sign expression as
\[
 \mathscr H_{m,q}
 =(2m+q-11)s^2+\frac12R\bigl(4mv+4(q+1)u-16R\bigr)s-12RD.
\]
At \((m,q)=(4,4)\), direct substitution gives
\begin{equation}\label{eq:Clifford-H44}
 \mathscr H_{4,4}
 =\frac{T(8u^2+T)}{16}+(u^2-T)v(2u+3v)\geq0,
\end{equation}
by \eqref{eq:box-constraints} and \(u,v\geq0\).
For every \(m\geq4\), \(q\geq4\),
\begin{align*}
 \mathscr H_{m,q}
 &=\mathscr H_{4,4}
   +\bigl(2(m-4)+q-4\bigr)s^2\\
 &\quad+2R\bigl((m-4)v+(q-4)u\bigr)s\geq0,
\end{align*}
and \(\kappa\geq16R\). Hence \cref{thm:cubic-unified-sextic}
applies in this range as well.
Finally,
\[
 f\sqrt{Rs}=\frac14F\sqrt{RS}.
\]
This positive factor preserves the normalized gradient and the
strict divergence sign, giving the result for the original \(F\).
\end{proof}

In the normalized variables of \eqref{eq:Clifford-normalized-data},
the sextic is generally independent of the cubic:
\[
 D-f^2=\frac{vT-F^2}{4}\geq0.
\]
This Cauchy--Schwarz defect need not vanish, so the polynomial
gradient-square criterion requiring \(D=f^2\) does not apply in general.

\paragraph*{Dimensions in which Clifford systems exist}

For a symmetric Clifford system $(A_0,\ldots,A_q)$ on $\R^{2m}$,
the Hurwitz--Radon admissibility condition
\cite[Chapter~0, pp.~3--4]{Shapiro2000} is \(q\leq\rho(m)\),
with the Clifford-system convention of \cite{FKM81,TkachevClifford10}.
Here, if $m=2^{4a+b}c$ with $c$ odd and $0\leq b\leq3$, then \(\rho(m)=8a+2^b\).
In particular, \(\rho(4)=4,\; \rho(5)=1,\; \rho(6)=2,\; \rho(7)=1\).
Thus an admissible pair with $q\geq3$ has either $m=4$ or $m\geq8$.
If $m=4$, then $q\in\{3,4\}$.

\begin{proof}[Sufficiency in \cref{thm:clifford-classification}]
By admissibility, every pair in
\eqref{eq:clifford-classification-list} has either \(m\geq8\),
\(q\geq3\), or \((m,q)=(4,4)\). Both cases are covered by
\cref{prop:Clifford-common-potential}.
\end{proof}

\subsection{Singular strata and tangent cones}
\label{sec:tangents}

From \eqref{eq:U}, the critical set is
\[
 \Sigma:=\{\nabla F=0\}
 =\Sigma_0\cup\Sigma_+,
\]
where
\[
 \Sigma_0=\{(0,z):z\in\R^{q+1}\},\qquad
 \Sigma_+=\{(y,0):\tau(y)=0\}.
\]
Indeed, $A_z^2=v\Id$ shows that $A_zy=0$ forces $z=0$ or $y=0$.

At a nonzero point of $\Sigma_+$, the gradients
$\nabla\tau_i=2A_i y$ are mutually orthogonal, so this stratum is
smooth there and has ambient codimension $2q+2$.  The stratum
$\Sigma_0$ has ambient codimension $2m$.  In the sufficient range
$m\geq4$, $q\geq3$, and with $N=2m+q+1$,
\begin{equation}\label{eq:Sigma-thin}
 \cH^{N-1}(\Sigma\cap \mathcal K)=0
 \quad\text{for every compact }\mathcal K.
\end{equation}

Define the generalized Simons cone
\[
 \mathcal S_k
 =\{(a,b)\in\R^k\oplus\R^k:|a|=|b|\}.
\]
Let $E=\{F<0\}$.

\Needspace{5\baselineskip}
\begin{proposition}[Unique tangents with multiplicity one]
\label{prop:tangent-cones}
\leavevmode
\begin{enumerate}[label=\textup{(\roman*)}]
\item At a nonzero point $p=(0,z_0)$ of $\Sigma_0$, choose orthogonal
coordinates
\[
 \R^{2m+q+1}=\R^m\oplus\R^m\oplus\R^{q+1}.
\]
The unique set blow-up is
\begin{equation}\label{eq:tangent-Sigma0}
 E_p^{(0)}=\{(a,b,\eta):|a|<|b|\},\qquad
 \partial E_p^{(0)}=\mathcal S_m\times\R^{q+1}.
\end{equation}
\item At $p=(y_0,0)\in\Sigma_+\setminus\{0\}$, after an orthogonal
identification
$\R^{2m+q+1}=\R^{q+1}\oplus\R^{q+1}
 \oplus\R^{2m-q-1}$, the unique set blow-up is
\begin{equation}\label{eq:tangent-Sigmaplus}
 E_p^{(+)}=\{(x,x',\xi^\perp):|x|<|x'|\},\qquad
 \partial E_p^{(+)}
 =\mathcal S_{q+1}\times\R^{2m-q-1}.
\end{equation}
\item At the vertex, every dilation fixes $E$ and $\cC_{\cA}$, so the
tangent cone is $\cC_{\cA}$ itself.
\end{enumerate}
In each case, the oriented current tangent is the boundary of the
stated set, and the varifold tangent has multiplicity one on its regular
part.  The perimeter measures converge under dilation.  These statements
hold for every admissible Clifford system, without an area-minimizing
assumption.
\end{proposition}

\begin{proof}
The Taylor argument in \cref{sec:tangent-instability-transfer} gives
the set tangents. We use exact local coordinates also to determine
the current and varifold limits. Near
$p=(0,z_0)$, $z_0\ne0$, the identity $A_z^2=|z|^2\Id$ gives
smooth positive and negative eigenspaces. Their dimensions are both
$m$, since $A_w$ interchanges them whenever $0\ne w\perp z$.
Choose smooth orthonormal frames and write $y=P(z)(a,b)$. Then
\[
 F(y,z)=|z|\bigl(|a|^2-|b|^2\bigr).
\]
The coordinate map has orthogonal derivative at $(0,z_0)$, since
its derivatives in $z$ are multiplied by $y=0$. This gives
\eqref{eq:tangent-Sigma0}, including the stated set side.

Near $p=(y_0,0)\in\Sigma_+$, $y_0\ne0$, the vectors
$e_i=A_i y_0/|y_0|$, $0\le i\le q$, are orthonormal.
Let $P_\perp$ be the orthogonal projection onto their common
orthogonal complement. The map
\[
 y\longmapsto\left(
 a_i(y)=\frac{\tau_i(y)}{2|y_0|}\ (0\le i\le q),\quad
 \xi^\perp=P_\perp(y-y_0)\right)
\]
has orthogonal derivative at $y_0$, because
$da_i|_{y_0}(\xi)=\langle e_i,\xi\rangle$, and is a smooth
local coordinate map. With $z=\eta$ and
\[
 x=\frac{a+\eta}{\sqrt2},\qquad
 x'=\frac{a-\eta}{\sqrt2},
\]
we have exactly
\[
 F(y,\eta)=2|y_0|a\cdot\eta
          =|y_0|\bigl(|x|^2-|x'|^2\bigr).
\]
This proves \eqref{eq:tangent-Sigmaplus} with its stated side.

In either case, after the stated orthogonal identification, a smooth
local diffeomorphism \(\Phi\) takes the quadratic sublevel model onto
\(E\), with \(\Phi(0)=p\) and \(D\Phi(0)=\Id\).
The rescaled maps \(\Phi_r(w)=(\Phi(rw)-p)/r\) converge to the identity
in \(C^1\) on compact sets.  Each quadratic model has locally finite
perimeter and multiplicity one on its regular boundary.  The area
formula for \(\Phi_r\), whose tangential Jacobians converge uniformly
to one, gives convergence of the perimeter measures and the associated
varifolds.  The $C^1$ convergence also gives $L^1_{\rm loc}$ convergence of the
sublevel sets, whose model boundaries are Lebesgue null, and identifies
the oriented current limit with the boundary of the stated model.  At the vertex, all rescalings
agree exactly by homogeneity.
\end{proof}

\begin{remark}
All nonzero algebraic singular points are genuine geometric
singularities: neither
\eqref{eq:tangent-Sigma0} nor \eqref{eq:tangent-Sigmaplus} is a
hyperplane. Thus \(\Sing\cC_\cA=\Sigma_0\cup\Sigma_+\).
The transverse Simons cones are area minimizing exactly when
$m\geq4$ and $q+1\geq4$, respectively
\cite{BDGG69,DePhilippisPaolini}. For a transverse parameter at least
two, this is also exactly the stable range \cite[Section~6.1]{Simons68}; the parameter-one planar
crossing is stable on its regular part but is not minimizing. In particular,
for $(m,q)=(4,3)$ both strata have the same tangent model
$\mathcal S_4\times\R^4$.  The corresponding multiplicity-one current tangents are area minimizing;
at regular points the tangent is a multiplicity-one hyperplane.

For the Hurwitz cubics \(F(x,y,z)=\operatorname{Re}((xy)z)\), with
\(x,y,z\in\mathbb K\) and \(\mathbb K=\mathbb C,\mathbb H,\mathbb O\),
the singular set is the union of the three mutually orthogonal
coordinate subspaces:
\[
 \Sing\{F=0\}
 =(\mathbb K\oplus0\oplus0)\cup(0\oplus\mathbb K\oplus0)
   \cup(0\oplus0\oplus\mathbb K).
\]
Indeed, \(|\nabla F|^2=|x|^2|y|^2+|y|^2|z|^2+|z|^2|x|^2\).
The complex case gives the explicit cubic tangent in
\cref{sec:tangent-instability-transfer}.
\end{remark}

\begin{lemma}
\label{lem:Sigmaplus-nonempty}
If $m>q$, then $\Sigma_+\setminus\{0\}$ is nonempty.
\end{lemma}

\begin{proof}
After an orthogonal change of coordinates, the first two Clifford
matrices have the standard form
\begin{equation}\label{eq:standard-first-two}
 A_0=
 \begin{pmatrix}\Id&0\\0&-\Id\end{pmatrix},
 \qquad
 A_1=
 \begin{pmatrix}0&\Id\\\Id&0\end{pmatrix},
\end{equation}
and, for $i\geq2$,
\begin{equation}\label{eq:standard-rest}
 A_i=
 \begin{pmatrix}0&E_i\\-E_i&0\end{pmatrix},
\end{equation}
where the $E_i$ are skew-symmetric orthogonal matrices.  Fix
$b\neq0$ in $\R^m$.  The $q$ vectors
$b,E_2b,\ldots,E_qb$ are mutually orthogonal.  Since $m>q$, choose
$a\neq0$ of the same length as $b$ and orthogonal to their span.
For $y=(a,b)$, all components of $\tau(y)$ vanish.
\end{proof}

\begin{proof}[Necessity in \cref{thm:clifford-classification}]
Suppose first that $m\leq3$.  The stratum
$\Sigma_0\setminus\{0\}$ has dimension $q+1$ and consists of genuine
singular points by \cref{prop:tangent-cones}.  Its ambient codimension
is $2m<8$.  Federer's codimension-seven bound for the singular set of an
area-minimizing hypersurface \cite{Federer70} therefore excludes area
minimization.

Now suppose $q\leq2$ and $m\geq4$.  Then $m>q$, so
\cref{lem:Sigmaplus-nonempty} applies.  The smooth stratum
$\Sigma_+\setminus\{0\}$ has ambient codimension $2q+2<8$ and again
consists of genuine singular points.  The same dimension bound excludes area minimization in this case as well.
\end{proof}

\subsection{The exceptional pair \texorpdfstring{$(4,3)$}{(4,3)}}

\label{sec:instability}

\paragraph*{Second variation and the radial Hardy inequality}

Fix $(m,q)=(4,3)$ and write
\[
 \Lambda=\cC_{\cA}\cap\Sph^{11},\qquad
 \mathscr S=\Sing(\cC_{\cA})\cap\Sph^{11},\qquad
 \mathscr S_0=\Sigma_0\cap\Sph^{11},\qquad
 \mathscr S_+=\Sigma_+\cap\Sph^{11}.
\]
Thus $\dim\Lambda=10$, while the two singular strata $\mathscr S_0$ and
$\mathscr S_+$ each have dimension $3$.
Let \(\mathcal D(\Lambda)\) be the closure of
\(C_c^\infty(\Lambda_{\reg})\) in the norm with square
\(\int_{\Lambda_{\reg}}(|\nabla\psi|^2+(1+|A_\Lambda|^2)\psi^2)\).
The radial Hardy inequality in cone dimension \(n=11\) has optimal
constant \((n-2)^2/4=81/4\).  Separating the radial variable in the
second variation of area therefore gives the quadratic form
\[
 \mathcal B(\varphi,\psi)=\int_{\Lambda_{\reg}}
 \left(\langle\nabla_\Lambda\varphi,\nabla_\Lambda\psi\rangle
       -|A_\Lambda|^2\varphi\psi+\frac{81}{4}\varphi\psi\right).
\]
It remains to find \(\phi\in\mathcal D(\Lambda)\) with
\(\mathcal B(\phi,\phi)<0\).

For the present link, the smooth compact singular strata and their
nondegenerate transverse quadratic models give
\begin{equation}\label{eq:link-tube}
 \cH^{10}(N_r(\mathscr S)\cap\Lambda)\leq Cr^7,
 \qquad |A_\Lambda|\leq C\rho^{-1},\qquad
 \rho=\dist(\cdot,\mathscr S).
\end{equation}
Indeed, the local coordinates in \cref{prop:tangent-cones} identify
the transverse zero sets with seven-dimensional quadratic cones in
eight normal variables.  Their volume element has radial order
\(\rho^6\dd\rho\), and their second fundamental forms have size
\(O(\rho^{-1})\).  Restricting to the unit sphere removes the radial
direction along the singular stratum.  The coordinate maps and their
inverses have bounded first and second derivatives on a finite cover
of the compact singular strata.  The same estimates therefore hold on
\(\Lambda\), which proves \eqref{eq:link-tube}.

\paragraph*{Curvature identities on the link}

Use $u,v,S,J$ from \cref{sec:invariants}. On
$\Lambda$ one has $R=u+v=1$ and $F=0$.

\begin{lemma}[Link curvature]
\label{lem:link-curvature}
On $\Lambda_{\reg}$,
\begin{align}
 |A_\Lambda|^2&=\frac{32}{S}-\frac{J}{2S^2},
 \label{eq:link-curvature}\\
 |\nabla_\Lambda S|^2&=J-16S^2.
 \label{eq:link-gradS}
\end{align}
\end{lemma}

\begin{proof}
The curvature identity in the proof of
\cref{prop:harmonic-power-Jacobi} is
\(|A_C|^2=K/(2S)-J/(2S^2)\). Here \eqref{eq:DeltaS} gives
\(K=64\) at \((m,q)=(4,3)\), \(R=1\), proving
\eqref{eq:link-curvature}: the radial principal curvature is zero,
so \(|A_C|=|A_\Lambda|\) at radius one.
Moreover, \(Z\perp\nabla F\) on \(F=0\) by \eqref{eq:UZ}, and homogeneity gives
\(\langle Z,(y,z)\rangle=4S\).  Hence
\(\nabla_\Lambda S=Z-4S(y,z)\), proving
\eqref{eq:link-gradS}.
\end{proof}

\paragraph*{Integration along the Hopf fibers}

For $(4,3)$, the standard block form allows us to extend the Clifford
system by one generator.  In the notation of
\eqref{eq:standard-first-two}--\eqref{eq:standard-rest}, put
$E_4=E_2E_3$ and
\[
 A_4=\begin{pmatrix}0&E_4\\-E_4&0\end{pmatrix}.
\]
Then $(A_0,\ldots,A_4)$ is a Clifford system on $\R^8$.  For
$\xi\in\Sph^7$, the map
\begin{equation}\label{eq:hopf-map}
 \Pi(\xi)
 =\bigl(\langle A_0\xi,\xi\rangle,\ldots,
        \langle A_4\xi,\xi\rangle\bigr)
\end{equation}
is the quaternionic Hopf map $\Sph^7\to\Sph^4$; see, for example,
\cite{FKM81}.  To verify its measure normalization, write
\(p_i(\xi)=\langle A_i\xi,\xi\rangle\).  If
$\xi=(a,b)\in\R^4\oplus\R^4$, then
$b,E_2b,E_3b,E_4b$ are mutually orthogonal with the same length;
this proves $|\Pi(\xi)|=1$.  The Clifford relations also give
\begin{equation}\label{eq:hopf-horizontal}
 \nabla_{\Sph^7}p_i=2(A_i\xi-p_i\xi),\qquad
 \langle\nabla_{\Sph^7}p_i,\nabla_{\Sph^7}p_j\rangle
 =4(\delta_{ij}-p_ip_j).
\end{equation}
Thus \(D\Pi\) multiplies lengths by \(2\) on the horizontal space
\((\ker D\Pi)^\perp\), and its four-dimensional Jacobian is the
constant \(2^4\).
The quaternionic Hopf fibers are congruent round \(3\)-spheres.
Integration over the fibers gives, for every integrable function \(h\),
\[
 \int_{\Sph^7}h(\Pi(\xi))\dd\xi
 =c_H\int_{\Sph^4}h(p)\dd p,
 \qquad c_H>0.
\]

The Hopf map also describes the second singular stratum: $\mathscr S_+$ is the inverse image under $\Pi$ of the two poles
$(0,0,0,0,\pm1)\in\Sph^4$.  Hence $\mathscr S_+$ has two connected
components, each a Hopf fiber diffeomorphic to $S^3$.

Set
\[
 \theta=\frac{T}{u^2}\in[0,1].
\]
If $s=\langle A_4\xi,\xi\rangle$, then $\theta=1-s^2$.
The following lemma expresses integration on the link in terms of
$u$ and $\theta$.

\begin{lemma}[Integration on the link]
\label{lem:hopf-coarea}
There is a constant $c_0>0$ such that every integrable function
$h=h(u,\theta)$ satisfies
\begin{align}
 \int_{\Lambda_{\reg}}h\,\dd\cH^{10}
 &=c_0\int_0^1\int_0^1
     h(u,\theta)\,
     u^2(1-u)^{1/2}
     \left(\frac{\theta}{1-\theta}\right)^{1/2}
     S^{1/2}\,\dd\theta\,\dd u,
 \label{eq:hopf-coarea}\\
 S&=4u(1-u)+u^2\theta,                                  \notag\\
 J&=16\left[
       4u(1-u)+\bigl(u+4(1-u)\bigr)u^2\theta
     \right].
 \label{eq:Clifford-link-SJ}
\end{align}
\end{lemma}

\begin{proof}
Write $y=\sqrt u\,\xi$ and $z=\sqrt{1-u}\,\zeta$, with
$\xi\in\Sph^7$ and $\zeta\in\Sph^3$.  Spherical measure on $\Sph^{11}$
decomposes, up to a positive constant, as
\[
 u^3(1-u)\,\dd u\,\dd\xi\,\dd\zeta.
\]
If
$\bar\tau=(\langle A_0\xi,\xi\rangle,\ldots,
\langle A_3\xi,\xi\rangle)$, then
$|\bar\tau|=\sqrt\theta$ and
\[
 F=u\sqrt{1-u}\,\langle\zeta,\bar\tau\rangle.
\]
For fixed $0<u<1$ and $\xi$ with $\theta>0$, the equation $F=0$
defines an equatorial $\Sph^2$ in the $\zeta$-sphere.  Along this equator,
\begin{equation}\label{eq:z-coarea}
 |\nabla_{\Sph^3,\zeta}F|=u\sqrt{1-u}\sqrt\theta.
\end{equation}
The coordinate exceptional sets have zero $\cH^{10}$-measure in
$\Lambda$. Indeed, $u=0$ and $u=1$ have dimensions three and seven.
For $0<u<1$, the set $\theta=0$ has dimension $1+3+3=7$, since
$\xi$ lies in two Hopf fibers and $\zeta$ is unrestricted; the set
$\theta=1$ has dimension $1+6+2=9$, since $\xi$ lies above an
equatorial $\Sph^3\subset\Sph^4$ and $\zeta$ lies in an equatorial
$\Sph^2$. We may therefore work on $0<u<1$, $0<\theta<1$ without
changing the link integral.

On $\{F=0\}\cap\Sph^{11}$, Euler's identity gives
$|\nabla_{S^{11}}F|=|\nabla F|=\sqrt S$.
Applying the coarea formula first on $\Sph^{11}$ and then on the
$\zeta$-sphere gives, for a constant $c_1>0$,
\[
 \int_{\Lambda_{\reg}}h\dd\cH^{10}
 =c_1\int_0^1u^2(1-u)^{1/2}
       \int_{\Sph^7}h(u,\theta)\frac{S^{1/2}}{\sqrt\theta}
       \dd\xi\,\dd u.
\]

It remains to integrate in $\xi$.  The Hopf map
\eqref{eq:hopf-map} has congruent fibers and constant Jacobian.
For the last coordinate $s$ on $\Sph^4$, spherical measure has density
proportional to $(1-s^2)\dd s$.  Since $\theta=1-s^2$, the two
branches of $s$ contribute
\[
 \frac{\theta}{\sqrt{1-\theta}}\,\dd\theta.
\]
Substitution into the preceding integral proves
\eqref{eq:hopf-coarea}; \eqref{eq:Clifford-link-SJ} follows from
the formulas for $S$ and $J$ in \eqref{eq:S} with
$R=1$ and $T=u^2\theta$.
\end{proof}

\paragraph*{A test function proving instability}

For $a\in\R$, put $\phi_a=S^a$ on $\Lambda_{\reg}$.  Combining
\cref{lem:link-curvature,lem:hopf-coarea} gives, whenever the
integral is finite,
\begin{align}
 \frac1{c_0}\mathcal B(\phi_a,\phi_b)
 &=\int_0^1\int_0^1
    u^2(1-u)^{1/2}
    \left(\frac{\theta}{1-\theta}\right)^{1/2}
    \Bigg[
       \left(ab+\frac12\right)JS^{a+b-3/2}              \notag\\
 &\hspace{27mm}
       -32S^{a+b-1/2}
       +\left(\frac{81}{4}-16ab\right)S^{a+b+1/2}
    \Bigg]\dd\theta\,\dd u.
 \label{eq:monomial-link-form}
\end{align}
For $a,b\in\{3/4,7/4\}$, every power of $S$ in
\eqref{eq:monomial-link-form} is a nonnegative integer.  Expanding
\eqref{eq:Clifford-link-SJ} and using only
\begin{align*}
 \int_0^1u^{i+2}(1-u)^{1/2}\dd u
   &=\mathrm B\left(i+3,\frac32\right),\\
 \int_0^1\theta^{j+1/2}(1-\theta)^{-1/2}\dd\theta
   &=\mathrm B\left(j+\frac32,\frac12\right),
\end{align*}
we obtain
\[
 \frac1{c_0\pi}
 \begin{pmatrix}
  \mathcal B(\phi_{3/4},\phi_{3/4})
   &\mathcal B(\phi_{3/4},\phi_{7/4})\\
  \mathcal B(\phi_{7/4},\phi_{3/4})
   &\mathcal B(\phi_{7/4},\phi_{7/4})
 \end{pmatrix}
 =
 \begin{pmatrix}
  \dfrac{96}{715}
   &\dfrac{471936}{1616615}\\[2mm]
  \dfrac{471936}{1616615}
   &\dfrac{571904}{969969}
 \end{pmatrix}.
\]
In particular, for
\[
 \phi=\phi_{3/4}-\frac12\phi_{7/4}
      =S^{3/4}\left(1-\frac12S\right),
\]
\begin{equation}\label{eq:negative-link-form}
 \mathcal B(\phi,\phi)
 =-c_0\pi\,\frac{9952}{969969}<0.
\end{equation}

Near either singular stratum of $\mathscr S$, the nondegenerate transverse
quadratic models give
\[
 S\asymp\rho^2,
 \qquad
 J=O(\rho^2),
 \qquad
 |\nabla_\Lambda S|=O(\rho),
\]
where $\rho=\dist(\cdot,\mathscr S)$; the last estimate also follows from
\eqref{eq:link-gradS}.  Therefore
\[
 |\phi|=O(\rho^{3/2}),
 \qquad
 |\nabla_\Lambda\phi|=O(\rho^{1/2}).
\]
In particular, both $|\nabla_\Lambda\phi|^2$ and
$|A_\Lambda|^2\phi^2$ are $O(\rho)$.  Since the transverse volume
element has order $\rho^6\dd\rho$, all terms in the form energy are
locally integrable.  Choose smooth cutoffs \(\chi_\eps\) that vanish
for \(\rho\leq\eps\), equal one for \(\rho\geq2\eps\), and satisfy
\(|\nabla_\Lambda\chi_\eps|\leq C/\eps\).
The estimate \eqref{eq:link-tube} and the preceding pointwise bounds give
\[
 \|(1-\chi_\eps)\phi\|_{\mathcal D}^2
 \leq C\eps^8\longrightarrow0.
\]
Indeed, the gradient and curvature terms are bounded by
\(C\int_0^{2\eps}\rho\,\rho^6\dd\rho=O(\eps^8)\);
the cutoff derivative contributes at most
\(C\eps^{-2}\eps^3\eps^7=O(\eps^8)\).
Thus \(\chi_\eps\phi\in C_c^\infty(\Lambda_{\reg})\) approximates
\(\phi\) in \(\mathcal D(\Lambda)\).
The cone stability formula \eqref{eq:exact-Hardy-gap} and
\eqref{eq:negative-link-form} prove
\cref{thm:Clifford-43-instability}.  Together with the sufficient and
dimension-obstruction arguments above, this also completes the proof
of \cref{thm:clifford-classification}.

\begin{remark}[Power multipliers at \((4,3)\)]
\label{sec:obstruction}
\label{prop:Clifford-power-obstruction}
Instability already excludes every positive smooth multiplier with the
two-sided sign near the regular cone: \cref{thm:Jacobi-trace} and
\eqref{eq:ground-state-identity} would otherwise imply stability. In particular,
no choice of exponents in
\[
 G_{\alpha,\gamma}=R^\alpha S^\gamma F,
 \qquad(\alpha,\gamma)\in\R^2,
\]
can give a subcalibration for \((m,q)=(4,3)\).
\Cref{app:Clifford-power-obstruction} additionally locates explicit
points in the actual invariant image with \(R,S>0\), \(F\ne0\),
and \(\nabla G_{\alpha,\gamma}\ne0\), where
\(F\,\diver(\nabla G_{\alpha,\gamma}/|\nabla G_{\alpha,\gamma}|)<0\).
\end{remark}

\subsection{Unequal dimensions: Hopf incidence cubics}

Let $d\in\{2,4,8\}$ and let $h_d:\R^{2d}\to\R^{d+1}$ be the
standard quadratic Hopf map, with $|h_d(y)|^2=|y|^4$. Set
\[
 F_d(x,y)=\langle x,h_d(y)\rangle,\qquad
 x\in\R^{d+1},\quad y\in\R^{2d},\quad u=|y|^2,\quad v=|x|^2.
\]
The bidegree is $(1,2)$, and the Hopf identities give
\begin{equation}\label{eq:incidence-closure}
 (\Gamma(I_i,I_j))_{I=(u,v,F_d)}=
 \begin{pmatrix}4u&0&4F_d\\0&4v&2F_d\\
 4F_d&2F_d&u^2+4uv\end{pmatrix},\qquad
 \Delta I=(4d,2(d+1),0).
\end{equation}
Its determinant is $16(u+4v)(u^2v-F_d^2)$, positive on an open set.
Thus the algebra has three independent generators of degrees $(2,2,3)$.
Furthermore $\cL(F_d)=-8(u+v)F_d$, and the singular set of the
minimal zero hypercone is $\{y=0\}$. At a nonzero point of this
stratum the transverse quadratic form has signature $(d,d)$.

These are the Clifford cubics with parameters $(m,q)=(d,d)$.
By \cref{thm:clifford-classification}, the cases $d=4,8$ are area
minimizing, in ambient dimensions $13,25$. For $d=2$, the transverse
Lawson cone in $\R^4$ is unstable, and the localization argument
following \eqref{bh:Lawson-margin} gives instability in $\R^7$.
For \(d=4\), the degree-six potential is critical:
\(6=(13-1)/2\). Substituting \(T=u^2\) into
\eqref{eq:Clifford-H44} and \eqref{eq:cubic-H-Jacobi-meaning} gives
\[
 \beta_G=\frac{9u^2}{(u+4v)^2}>0
 \quad\text{on }C_{\reg},\qquad
 \inf_{C_{\reg}\cap\Sph^{12}}\beta_G=0.
\]
As \(u\to0\) near the singular stratum, the bound degenerates.
It proves stability but gives no uniform positive Hardy gap.
The three-generator description uses $|h_d(y)|^2=u^2$.
For a general Clifford system, $|\tau(y)|^2$ is a separate quartic
invariant. The computations in \cref{sec:invariants} use the
four-generator algebra $\R[u,v,|\tau|^2,F]$.

The corresponding regular levels on
\(\Sph^d\times\Sph^{2d-1}\) are isoparametric hypersurfaces.
Indeed, on the product of unit round spheres
\(P_d=\Sph^d\times\Sph^{2d-1}\), put \(f_d=F_d|_{P_d}\).
The partial degrees \((1,2)\) and \eqref{eq:incidence-closure} give
\[
 |\nabla_{P_d}f_d|^2=5(1-f_d^2),\qquad
 \Delta_{P_d}f_d=-5d f_d.
\]
Thus the levels \(f_d=t\), \(-1<t<1\), are isoparametric,
and \(f_d=0\) is minimal in \(P_d\).
